% Gate diode and gate leakage (jfet-gate): the set-up a person would wire, semantically the deck's.
%   VGF gf gsw 0     a 0 V source from the swept node into copy A's gate   -> ammeter I_G (measured)
%   JF 0 gf 0 DUT    copy A, drain and source grounded                     -> njfet
%   VGSW gsw 0, .dc  the gate drive, swept forward (and reverse)           -> source V_GS (set)
%   VF0 sw df0 0     a 0 V source from the swept node to copy B's drain    -> ammeter I_D (measured)
%   VGB0 gb0 0 -1.5  copy B's gate bias; its current is IG                 -> source V_GS (set) + ammeter I_G (measured)
%   JD0 df0 gb0 0    copy B, source grounded                               -> njfet
%   VSW sw 0, .dc    the swept drain supply                                -> source V_DS (set)
%   Three decks: forward and reverse are copy A with two sweeps; the field deck is copy B, four copies.
\documentclass[tikz, border=10pt]{standalone}
\usepackage{bench}
\begin{document}
\begin{tikzpicture}
\benchtitle{0,7.0}{Gate diode and gate leakage}
\def\gnd{1.0}
% the transistor
\node[njfet] (QA) at (4,3) {};
\node[anchor=west, font=\small, align=left] at ($(QA.D)+(0.15,-0.55)$) {A};
% source to ground
\draw (QA.S) -- (QA.S |- 0,\gnd) node[ground]{};
\nodename{anchor=west}{$(QA.S)+(0.1,-0.3)$}{0}
% the transistor
\node[njfet] (QB) at (11,3) {};
\node[anchor=west, font=\small, align=left] at ($(QB.D)+(0.15,-0.55)$) {B};
% source to ground
\draw (QB.S) -- (QB.S |- 0,\gnd) node[ground]{};
\nodename{anchor=west}{$(QB.S)+(0.1,-0.3)$}{0}
% the drain grounded
\draw (QA.D) -- ++(0,0.5) node[ground, rotate=180]{};
\nodename{anchor=west}{$(QA.D)+(0.1,0.25)$}{0}
% the gate: a voltage source from ground, through an ammeter
\coordinate (Ag) at (QA.G);
\coordinate (Agam) at ($(Ag)+(-1.7,0)$);
\ammeterh{Agam}{Ag}{$I_G$}
\coordinate (Agsrc) at ($(Agam)+(-1.4,0)$);
\draw (Agam) -- (Agsrc);
\draw[set] (Agsrc) to[V, l_={$V_{GS}$}] (Agsrc |- 0,\gnd) node[ground, black]{};
\nodename{anchor=south}{$(Agam)+(-0.4,0.02)$}{gf}
% the ammeter in the drain, up to the rail
\coordinate (Btop) at ($(QB.D)+(0,1.7)$);
\ammeter{QB.D}{Btop}{$I_D$}
\nodename{anchor=west}{$(QB.D)+(0.1,0.25)$}{df0}
% the drain supply at the right, + up, - to ground
\coordinate (Bvtop) at ($(Btop)+(2.6,0)$);
\draw (Btop) -- (Bvtop);
\draw[set] (Bvtop) to[V, l={$V_{DS}$}] (Bvtop |- 0,\gnd) node[ground, black]{};
\nodename{anchor=south}{$(Btop)+(1.3,0.05)$}{sw}
% the gate: a voltage source from ground, through an ammeter
\coordinate (Bg) at (QB.G);
\coordinate (Bgam) at ($(Bg)+(-1.7,0)$);
\ammeterh{Bgam}{Bg}{$I_G$}
\coordinate (Bgsrc) at ($(Bgam)+(-1.4,0)$);
\draw (Bgam) -- (Bgsrc);
\draw[set] (Bgsrc) to[V, l_={$V_{GS}$}] (Bgsrc |- 0,\gnd) node[ground, black]{};
\nodename{anchor=south}{$(Bgam)+(-0.4,0.02)$}{gb0}
% the definitions
\definitions{0,-0.1}{%
  A: $V_{GS}$ swept 0 to $+0.9$ V in 5 mV steps (forward), and 0 to $-40$ V in 0.1 V steps (reverse), drain and source grounded;\quad read $I_G$\\[2pt]
  B: $V_{GS}$ the bias for $I_D = 100$ \textmu A, 1, 3 and 10 mA (from the transfer sweep), one copy each;\quad $V_{DS}$ swept 0 to 25 V in 0.1 V steps\\[2pt]
  read $I_G$ against $V_{DG} = V_{DS} - V_{GS}$;\quad $I_{GSS}$ at 10, 20 and 30 V from the reverse sweep%ENDDEFS
}
\end{tikzpicture}
\end{document}
